% =====================================================================
%  Soil Health Prediction from IoT Sensor Data Using Machine Learning
%  in Sub-Saharan Africa: A Systematic Literature Review and
%  Methodological Framework
%
%  Adrian Rugaba Asiimwe -- S24B23/054
%  Department of Computer Science, Uganda Christian University
%
%  REVISED SOURCE. Every change from the reviewed version is marked
%  with a numbered comment, IMP1 ... IMP18, so the two can be diffed.
%
%  Anything marked \TBC{} must be filled from the screening log before
%  submission. None of those values may be inferred or estimated:
%  they are counts and extractions that only the reviewer's records
%  can supply.
%
%  Figures: prisma.pdf and pipeline.pdf must be placed beside this file.
% =====================================================================
\documentclass[conference]{IEEEtran}
\usepackage{amsmath,amssymb}
\usepackage{graphicx}
\usepackage{booktabs}
\usepackage{array}
\usepackage{url}

% Visible marker for values that must still be extracted.
\newcommand{\TBC}{\textbf{[TBC]}}
\renewcommand{\arraystretch}{1.15}

\begin{document}

% IMP18: title scoped. The first half describes the reviewed literature, where
% predicting soil properties from IoT sensors is what the studies actually do.
% The second half now names what this paper contributes, so the sensor claim no
% longer transfers onto the framework. See the title alternatives in the notes.
\title{Soil Health Prediction from IoT Sensor Data Using Machine Learning
in Sub-Saharan Africa: A Systematic Literature Review and a Framework for
Validation and Sensor-Robustness Testing}

\author{\IEEEauthorblockN{Adrian Rugaba Asiimwe}
\IEEEauthorblockA{Department of Computer Science\\
Uganda Christian University, Mukono, Uganda\\
Reg no: S24B23/054}}

\maketitle

% IMP1: the abstract now states the evidence status of the framework.
% The reviewed version described the framework as if it were pending;
% a reference implementation exists, so say so and label its data.
\begin{abstract}
Soil degradation threatens agricultural productivity and food security, and
many farmers in sub-Saharan Africa (SSA) have little access to laboratory
soil testing, which is slow and requires specialised equipment. Internet of
Things (IoT) field sensors combined with machine learning (ML) are a possible
alternative: sensors measure soil properties such as moisture and pH in the
field, and ML models estimate laboratory values from the sensor readings. This
systematic literature review (SLR) examines how ML has been used to predict
soil properties from field sensor data, how the models were validated, and how
robust they are to noisy sensor inputs. It follows PRISMA 2020 and the
guidelines of Kitchenham and Charters. Scopus and Google Scholar were searched,
which returned 262 records; after duplicate removal, screening, full-text
assessment and backward snowballing, 20 studies were included. No study that
predicts soil properties from field sensor data with ML on Ugandan soils was
found in the sources searched. The paper then sets out a methodological
framework for an empirical study, comprising a sensor simulation and an
analysis of public soil datasets, and reports a reference implementation of
that framework. The implementation uses laboratory measurements as proxies and
a simulated sensor pathway; no part of it yet ingests field sensor data. Its
results are currently obtained from a labelled synthetic stand-in dataset and
are reported as such.
\end{abstract}

\begin{IEEEkeywords}
soil health, IoT sensors, machine learning, validation design, sensor
simulation, spatially blocked cross-validation, precision agriculture,
sub-Saharan Africa, Uganda
\end{IEEEkeywords}

% =====================================================================
\section{Introduction}
% =====================================================================
Soil degradation affects an estimated 1,036 to 1,470 million hectares worldwide
and threatens long-term agricultural productivity and food security~\cite{ref1}.
The range is wide because published estimates define degradation and assess it
in different ways. Uganda's soils are reported to be highly degraded, low in
nitrogen, phosphorus and potassium, and generally acidic~\cite{ref2}. Soil
conditions vary between regions, so this finding applies to the areas covered
by~\cite{ref2} and not necessarily to all of Uganda.

Farmers can respond only if they know what their soil lacks, yet limited access
to soil testing services has left many farmers in districts such as
Nakapiripirit and Amudat struggling with declining fertility~\cite{ref3}.
Without testing, farmers judge fertility by vegetation cover, weeds, soil colour
and previous yield. Elaete et al.\ compared these four indicators with Palin
test-kit results for 461 beneficiaries of an input-subsidy project in
Uganda~\cite{ref2}. Depending on the indicator, between 61 percent (previous
yield) and 71 percent (vegetation cover) of farmers misclassified their soil
fertility.

Laboratory analysis remains the scientific standard for soil health assessment,
but it is slow and out of reach for many farmers. Typical procedures include
atomic absorption spectrometry, the Kjeldahl method for total nitrogen, and
spectrophotometric analysis for boron and phosphorus~\cite{ref1}. All require
specialised equipment, trained personnel and controlled conditions that are
seldom available outside major research institutions. Visual indicators are
unreliable and laboratory tests are hard to reach, so farmers need an automated,
field-based method that returns nutrient-specific estimates at many points on a
farm.

Field sensors combined with ML offer such a method. Sensors measure soil
moisture, pH and nutrient levels in real time~\cite{ref4}. Moisture determines
water availability, and pH controls how much of each nutrient plants can
absorb~\cite{ref5}. ML models trained on paired sensor and laboratory data can
estimate nutrient levels directly from sensor readings; the properties predicted
include nitrogen, phosphorus and potassium~\cite{ref4,ref6}. Reviews identify
soil sensors as the most frequently employed technology for soil health
assessment~\cite{ref7}, and portable spectral analysis with edge artificial
intelligence has yielded near-laboratory accuracy for in-situ
measurement~\cite{ref10,ref11}.

\noindent\textbf{Farming context.} About 7 million Ugandan households, 80 percent
of the total, cultivated land or reared livestock in 2019, with an average
holding of 1.3~ha~\cite{ref15}. The 2024 census reports that 17 percent of
households engage in commercial agriculture, from 27 percent in Bunyoro and
Buganda to 7 percent in Karamoja~\cite{ref16}. The sources found do not describe
the crops or soil management of commercial farms; the 2008/09 Census of
Agriculture used a separate form for large-scale and institutional
farms~\cite{ref17}. The review therefore records farm type (smallholder or
commercial) for each included study.

\subsection{Problem Statement and Research Gap}
Ugandan research has compared farmers' indicators with test-kit results and
called for portable test kits that can be used
locally~\cite{ref2}, but it has not built or evaluated a predictive model.
% IMP2: "affordable" removed; the call in [2] is for portable, locally usable
% kits, and the review does not frame the problem economically.
Sensor-based ML studies exist elsewhere~\cite{ref4,ref6,ref11}, but the
literature reviewed here leaves three questions insufficiently examined:
(i) whether such models work on Ugandan soils; (ii) how the models were
validated, since random cross-validation on spatially clustered samples can look
more accurate than it will be on new fields; and (iii) how robust they are to the
noise expected in field sensor data. Section~\ref{sec:gaps} states each gap from
the reviewed studies.

\subsection{Aims and Research Questions}
The review is guided by three research questions.

\begin{itemize}
\item \textbf{RQ1:} Which IoT or field sensors and soil parameters are used in
soil health monitoring, and how well do their readings agree with laboratory
values? \emph{Measured by:} sensor types and parameters per study, agreement
with laboratory values ($R^2$, RMSE or bias), and the deployment challenges
reported (connectivity, power, calibration and maintenance, skills, data
availability, usability).
\item \textbf{RQ2:} Which machine learning models are most commonly used to
predict soil properties from sensor data, and how do they perform in terms of
accuracy and data requirements? \emph{Measured by:} the number of included
studies using each model family (a study counts once per family), reported
$R^2$, RMSE or accuracy compared only between studies reporting the same metric
for a similar task, and dataset sizes.
\item \textbf{RQ3:} How were the ML models validated, and how robust are they to
noisy sensor inputs? \emph{Measured by:} the validation type (independent or
spatially blocked data, random split, or cross-validation only) and whether
input noise or sensor error was tested.
\end{itemize}

The objectives are to review field sensing of soil properties, identify the ML
models used to predict them, examine how those models were validated and tested
for noise, and identify the gaps that support a methodological framework for an
empirical study on Ugandan farms.

\subsection{Scope}
GIS mapping was considered during planning and removed from the scope. The
studies reviewed did not provide a basis for justifying maps as a component of
the prediction problem, and mapping adds interpolation choices that are separate
from the ML question. Spatial information is used here only to design validation,
by holding out whole areas.
% IMP3: the second clause of the original sentence named cost; the exclusion of
% economic analysis is kept, the wording is not.
Economic analysis is outside the scope of the review.

% =====================================================================
\section{Method}
% =====================================================================
\subsection{Research Approach and Protocol}
This study is an SLR. A systematic approach makes the search, selection and
analysis of studies explicit and repeatable, which reduces selection bias
compared with a narrative review. The review was planned and is reported
following the PRISMA 2020 statement~\cite{ref12}, supported by the SLR
guidelines of Kitchenham and Charters~\cite{ref13} and the PRISMA-S extension
for reporting searches~\cite{ref14}. The protocol sets out the research
questions, databases, search strings, criteria, quality assessment items and data
extraction form. The review proceeds in four stages: identification, screening,
eligibility assessment and final inclusion (Fig.~\ref{fig:prisma}).

\noindent\textbf{Operational definitions.} ``Farmers'' means people who grow crops
or manage agricultural land, whatever the size of their holding.

\noindent\textbf{Sensor.} A sensor is a field-deployable or proximal device that
measures a soil property without laboratory processing (for example
electrochemical, capacitive, optical, spectral or laser-induced breakdown
spectroscopy devices). Laboratory instruments and airborne or satellite imagery
are not counted.

\noindent\emph{Scope note.} This definition governs which studies are included.
The framework in Section~\ref{sec:method} models a sensor of this kind in
simulation; it does not deploy one, and this paper reports no field sensor data.

\noindent\textbf{Evidence origin.} Each included study is labelled
\emph{SSA-demonstrated} (data collected in SSA) or \emph{transferable} (data from
another region, with the authors stating that the method applies to SSA).
Findings are reported separately for the two labels.

\noindent\textbf{Validation types.} \emph{Independent} validation tests on data
from separate fields, sites or seasons. \emph{Spatially blocked} validation holds
out whole areas. \emph{Random split or cross-validation} divides samples at
random, so near neighbours can appear in both training and test sets.

\subsection{Databases Searched}
Two sources were searched (Table~\ref{tab:databases}).

\begin{table}[t]
\caption{Databases searched and records identified.}
\label{tab:databases}
\centering
\footnotesize
\begin{tabular}{p{1.05cm}p{4.4cm}l}
\toprule
Database & Fields searched & Records\\
\midrule
Scopus & Title, abstract, keywords & 112\\
Google Scholar & Title and full text; up to the first 100 results per query, sorted by relevance & 150\\
\midrule
Total & & 262\\
\bottomrule
\end{tabular}
\end{table}

Scopus was chosen as the main source for its multidisciplinary index with strong
agricultural, environmental and engineering coverage. Google Scholar was added as
a supplementary source for studies missed by Scopus, including regional work. Web
of Science and ScienceDirect were not searched; this is acknowledged as a
limitation (Section~\ref{sec:limits}). Google Scholar results change over time and
cannot be exported directly, so they were retrieved with Publish or Perish. Each
query was sorted by relevance and limited to the first 100 results, patents were
excluded, and the pooled results of all queries gave 150 records after
citation-only entries were removed.

\subsection{Search Keywords and Search Strings}
Keywords were derived from the research questions and fall into three concept
groups: soil, technology and context. Synonyms within a group are joined with OR,
and the three groups with AND. The master Boolean string is:

\begin{quote}
\footnotesize\ttfamily\raggedright
("soil health" OR "soil fertility" OR "soil quality" OR "soil nutrient*" OR
"soil moisture" OR "soil pH" OR "soil organic carbon") AND ("machine learning"
OR "deep learning" OR "predictive model*" OR "IoT" OR "Internet of Things" OR
"wireless sensor*" OR "soil sensor*" OR "spectroscopy" OR "edge AI") AND
("farmer*" OR "farming" OR "precision agriculture" OR "sub-Saharan Africa" OR
"East Africa" OR "Africa" OR "Uganda" OR "Kenya" OR "Tanzania" OR "Ethiopia" OR
"Rwanda" OR "Nigeria" OR "Ghana")
\end{quote}

The technology group deliberately joins the sensor and ML terms with OR so that
studies using only one of them are retrieved; the requirement of both (IC5) is
applied during screening. Remote sensing is not used as a search term, because
the review concerns sensors that farmers or extension workers can operate in the
field. The context group contains general farming terms, so studies from other
regions are retrieved and then judged for applicability to SSA under IC6. In
Scopus the string is wrapped in \texttt{TITLE-ABS-KEY(\ldots)}. Google Scholar
does not support long Boolean strings, so the string was split into shorter
queries, each combining one term from each concept group, and the results were
pooled.

\subsection{Inclusion and Exclusion Criteria}
A study must meet all inclusion criteria (Table~\ref{tab:include}) and none of
the exclusion criteria (Table~\ref{tab:exclude}).

\begin{table}[t]
\caption{Inclusion criteria.}
\label{tab:include}
\centering
\footnotesize
\begin{tabular}{@{}p{0.75cm}p{6.9cm}@{}}
\toprule
ID & Inclusion criterion\\
\midrule
IC1 & Peer-reviewed journal article or conference paper\\
IC2 & Published from January 2015 onward\\
IC3 & Written in English\\
IC4 & Addresses soil health, fertility or quality (e.g.\ nutrients, pH, moisture, organic matter)\\
IC5 & Uses an IoT or field sensor together with ML to predict or estimate a soil property\\
IC6 & Set in SSA; or set in another tropical or low-resource region \emph{and} the authors explicitly discuss transfer to SSA\\
IC7 & Reports enough technical detail (sensors, models, data) to allow data extraction\\
\bottomrule
\end{tabular}
\end{table}

\begin{table}[t]
\caption{Exclusion criteria.}
\label{tab:exclude}
\centering
\footnotesize
\begin{tabular}{@{}p{0.75cm}p{6.9cm}@{}}
\toprule
ID & Exclusion criterion\\
\midrule
EC1 & Not peer-reviewed (e.g.\ theses, editorials, white papers, posters)\\
EC2 & Published before 2015 or not in English\\
EC3 & Focuses on topics other than soil, such as crop disease, weather or livestock only\\
EC4 & Does not combine sensor data with ML (fails IC5)\\
EC5 & Opinion or position paper with no technical or empirical content\\
EC6 & Duplicate or earlier version of a study already included\\
EC7 & Full text not accessible\\
EC8 & Outside SSA and not applicable to it (fails IC6)\\
EC9 & Insufficient technical detail for data extraction (fails IC7)\\
\bottomrule
\end{tabular}
\end{table}

The start date of 2015 focuses the review on the period in which low-power IoT
hardware and applied ML became widely used in agriculture. IC6 admits studies
from other regions only when the authors discuss transfer to SSA, not merely
because they mention smallholders or the tropics.

\begin{table}[t]
\caption{Data extraction fields.}
\label{tab:fields}
\centering
\footnotesize
\begin{tabular}{@{}p{1.5cm}p{4.9cm}c@{}}
\toprule
Category & Data extracted & RQ\\
\midrule
Study details & Authors, year, title, venue, country or region, crop or farming system, scale, farm type, evidence origin & All\\
Sensing & Sensor types, soil parameters, power source, connectivity, calibration, deployment setting, agreement with laboratory values & RQ1\\
Deployment & Reported challenges (connectivity, power, calibration and maintenance, skills, data availability, usability) and proposed solutions & RQ1\\
Machine learning & Algorithm(s) and model family, input features, dataset size, metrics (accuracy, $R^2$, RMSE) & RQ2\\
Validation & Validation type (independent set, spatial blocks, random split or cross-validation), and whether robustness to input noise was tested & RQ3\\
Findings and limitations & Main results, stated limitations and gaps, quality score & All\\
\bottomrule
\end{tabular}
\end{table}

\subsection{Screening, Selection and Duplicate Removal}
Selection followed six stages: (1) identification, in which all records from the
database searches were exported and combined; (2) duplicate removal; (3) title
and abstract screening, where clearly failing records were excluded and uncertain
ones moved on; (4) full-text assessment, where the reason for each exclusion was
recorded with codes EC1 to EC9; (5) backward snowballing of the reference lists
of included studies, with the same screening steps; and (6) final inclusion of
studies from databases and snowballing for quality assessment, extraction and
synthesis.

Screening decisions and exclusion reasons were recorded in a screening log.
Records were exported in RIS or BibTeX format into a reference manager.
Duplicates were detected first by matching DOIs, then by matching title, first
author and publication year. Possible matches were checked manually before
deletion, and the version with the most complete metadata was kept.

\noindent\textbf{Single-reviewer screening.} The author screened all records, so no
second reviewer was used and inter-rater agreement was not measured. This is the
main source of selection bias in this review, because one person's reading of the
criteria decided every inclusion and exclusion. Every decision and its reason was
logged so that a second reviewer can check it.

\subsection{Data Extraction}
Data were recorded in a standardised spreadsheet form, which was first tested on
a small sample of included studies and then refined before full extraction.
Fields were chosen so that every research question can be answered
(Table~\ref{tab:fields}). Information a study does not report is recorded as
``not reported'' instead of being estimated.

\subsection{Quality Assessment}
Each study receives a quality score out of 6 using the checklist in
Table~\ref{tab:quality}.

\begin{table}[t]
\caption{Quality assessment checklist (maximum 6 points).}
\label{tab:quality}
\centering
\footnotesize
\begin{tabular}{@{}p{0.6cm}p{3.1cm}p{3.8cm}@{}}
\toprule
Item & Question & Scoring\\
\midrule
Q1 & Are the aim and farming context clearly stated? & 1 yes; 0 no\\
Q2 & Are the data source and sample size described? & 1 both; 0.5 one; 0 neither\\
Q3 & Is the method detailed enough to repeat? & 1 yes; 0.5 partly; 0 no\\
Q4 & Was the result validated on independent data? & 1 independent set or spatially blocked validation; 0.5 random split or $k$-fold CV only; 0 none\\
Q5 & Are performance metrics reported? & 1 yes; 0.5 partly; 0 no\\
Q6 & Are limitations discussed? & 1 yes; 0.5 briefly; 0 no\\
\bottomrule
\end{tabular}
\end{table}

The six items cover the main stages of a study (aim, data, method, validation,
results and reflection), in line with the general quality concerns discussed
in~\cite{ref13}. They carry equal weight because no agreed weights exist for
studies that mix sensing and ML, and unequal weights would add subjective
judgement; a sensitivity check (below) tests whether this choice affects the
findings.

Validation (Q4) distinguishes testing on independent data from
cross-validation. Random $k$-fold cross-validation reuses data from the same
samples and can overstate accuracy on new fields, so it earns half a point; an
independent dataset or spatially blocked validation earns a full point.

The score is not used to exclude studies. Studies scoring below 3 are retained
but (i) are marked in every table, (ii) are left out of the numeric ranges for
accuracy, and (iii) are cited in the synthesis only as supporting evidence beside
higher-scoring studies. Each finding is also reported with and without these
studies, so the effect of the weighting can be seen. Because studies are expected
to differ in sensors, models, datasets and metrics, no meta-analysis is planned,
and performance figures are compared only between studies reporting the same
metric for a similar task.

% =====================================================================
\section{Selection of Studies}\label{sec:selection}
% =====================================================================
The searches identified 262 records (Scopus 112, Google Scholar 150). After 31
duplicates were removed, 231 records were screened by title and abstract, and 174
were excluded. Of the 57 reports sought for retrieval, 6 could not be retrieved
(EC7) and 51 were assessed in full text. Full-text assessment excluded 34 reports
(EC1--EC2: 8; EC3: 9; EC4: 10; EC8: 4; EC9: 3), leaving 17 studies. Backward
snowballing of their reference lists added 3 studies, giving 20 studies for the
synthesis (Fig.~\ref{fig:prisma}).

\begin{figure}[t]
\centering
% Place the PRISMA diagram beside this file as prisma.pdf (or .png).
\includegraphics[width=\columnwidth]{prisma.pdf}
\caption{PRISMA 2020 flow diagram of the study selection process. Reasons for
full-text exclusion are given in Section~\ref{sec:selection}.}
\label{fig:prisma}
\end{figure}

\noindent\textbf{Search yield.} The string requires a term from each of three
concept groups, so the yield was limited to studies that mention soil, a
technology and a farming context together. Section~\ref{sec:limits} discusses the
effect of this on coverage.

% =====================================================================
\section{Included and Background Studies}
% =====================================================================
% IMP4: the reviewed version used one table for two different populations and
% then contradicted itself: Section IV called its rows background only, while
% Section VII counted "the seven primary studies in Table VI". The table is
% split here. Table VI keeps background and review sources, which are cited for
% context only. Table VIII holds the included studies, one row per study, and is
% the only source used for the counts in Section V.
Background studies are cited for context and are not part of the synthesis in
Section~\ref{sec:results}. A study is background only if it fails IC5, that is,
if it uses a field sensor or ML but not both. The seven primary background
studies are marked in Table~\ref{tab:background} together with the two reviews;
none of them covers field sensing, ML prediction and evaluation in SSA
together, whereas the proposed direction does.

\begin{table}[t]
\caption{Background and review sources. None of these rows contributes to the
counts in Section~\ref{sec:results}.}
\label{tab:background}
\centering
\footnotesize
\begin{tabular}{@{}p{1.55cm}p{3.2cm}ccc@{}}
\toprule
Study & Contribution / limitation & Sensor & ML & SSA\\
\midrule
Elaete et al.~\cite{ref2} & Farmer indicators vs Palin test-kit results in Uganda; no predictive model & No & No & Yes\\
Abdullah et al.~\cite{ref9} & IoT nutrient monitoring system, tested in Indonesia & Yes & No & No\\
Chlouveraki et al.~\cite{ref1} & Soil properties from hyperspectral imaging and ML & No & Yes & No\\
Erler et al.~\cite{ref10} & Mobile laser-induced breakdown spectroscopy; field-deployable, so a sensor under Section II-A & Yes & No & No\\
Mamabolo et al.~\cite{ref7} & Review of precision agriculture technologies & Review & Review & No\\
Getahun et al.~\cite{ref8} & Systematic review of precision agriculture technologies & Review & Review & No\\
\midrule
\multicolumn{5}{@{}p{7.4cm}@{}}{\emph{Included studies are listed in Table~\ref{tab:included}, not here.}}\\
\midrule
Proposed direction & Sensor-based ML prediction with spatial validation and noise tests on Ugandan soils & Yes & Yes & Yes\\
\bottomrule
\end{tabular}
\end{table}

% =====================================================================
\section{Results by Research Question}\label{sec:results}
% =====================================================================
% IMP5: this is the substantive change. In the reviewed version Sections V-A to
% V-C described how the studies would be grouped and then gave no counts; the
% 20 included studies appeared in no table. The three subsections below are now
% backed by Table VIII (one row per included study) and Table IX (the RQ3
% counts). Fill every \TBC from the extraction sheet. Do not estimate any of
% them, and do not write a count that Table VIII does not support.
All statements in this section are derived from Table~\ref{tab:included}, which
records each included study once. Statements marked ``Background'' come from
Table~\ref{tab:background} and are not findings of the review.

\subsection{RQ1: Sensors and Soil Parameters}
Table~\ref{tab:included} records the sensor type and soil parameter for each
study. Agreement with laboratory values is compared only between studies
reporting the same metric for the same property, and results are reported
separately for SSA-demonstrated and transferable studies. Deployment challenges
(connectivity, power, calibration and maintenance, skills, data availability,
usability) are counted across studies and reported in
Table~\ref{tab:rq3counts}.

\emph{Reported:} \TBC{} of the 20 included studies report agreement with
laboratory values; of these, \TBC{} are SSA-demonstrated. The most frequently
reported deployment challenge is \TBC{} (count \TBC{}).

\noindent\emph{Background.} Field sensors measuring soil moisture, pH and
nutrient levels in real time are the basis of IoT soil
monitoring~\cite{ref4}, and reviews report soil sensors as the most frequently
used technology for soil health assessment~\cite{ref7}.

\subsection{RQ2: Machine Learning Models}
Models are grouped into families (for example random forest, support vector
machine, neural network and regression), and a study counts once per family.
Accuracy is compared only between studies reporting the same metric for a similar
task, and studies scoring below 3 are left out of the ranges. Dataset sizes are
recorded because they limit which models can be trained.

\emph{Reported:} Table~\ref{tab:rq3counts} gives the number of included studies
per model family. The most frequently used family is \TBC{} (\TBC{} of 20
studies). Reported $R^2$ values range from \TBC{} to \TBC{} across the \TBC{}
studies reporting $R^2$ for a comparable task, and dataset sizes range from
\TBC{} to \TBC{} samples.

\noindent\emph{Background.} ML and deep learning models trained on paired sensor
and laboratory data estimate nutrient levels from sensor
readings~\cite{ref4,ref6}.

\subsection{RQ3: Validation and Robustness}
Each ML study is classed as independent, spatially blocked, or random-split or
cross-validation only, and checked for whether it tested robustness to input
noise or sensor error (Table~\ref{tab:included}, columns
\emph{Validation} and \emph{Noise}).

\emph{Reported:} of the 20 included studies, \TBC{} used an independent or
spatially blocked design, \TBC{} used a random split or cross-validation only,
and \TBC{} did not report a validation design that could be classified.
\TBC{} studies tested robustness to input noise. These counts are the central
evidence for the gap this review identifies: the share of reported accuracy that
rests on validation which can overstate performance on new fields is \TBC{} of
the included studies.

% ---------------------------------------------------------------------
% Table VIII: the included studies. One row per study, 20 rows.
% Column sources: Study, Region, Sensor, Model, n from Table IV fields
% "Study details", "Sensing", "Machine learning"; Validation and Noise from
% the "Validation" field; Q is the score from Table V (maximum 6).
% ---------------------------------------------------------------------
\begin{table*}[t]
\caption{Included studies ($n = 20$), extracted from the standardised form in
Table~\ref{tab:fields}. Validation: I = independent, SB = spatially blocked,
R = random split or cross-validation only, NR = not reported. Noise: whether
robustness to input noise or sensor error was tested. Q = quality score out of 6
(Table~\ref{tab:quality}); scores below 3 are marked $\dagger$ and excluded from
the accuracy ranges in Section~V-B.}
\label{tab:included}
\centering
\footnotesize
\begin{tabular*}{\textwidth}{@{\extracolsep{\fill}}cl p{2.5cm} p{3.0cm} p{1.8cm} r c c c@{}}
\toprule
\# & Study & Region & Sensor, parameters & Model family & $n$ & Valid. & Noise & Q\\
\midrule
1 & \TBC & \TBC & \TBC & \TBC & \TBC & \TBC & \TBC & \TBC\\
2 & \TBC & \TBC & \TBC & \TBC & \TBC & \TBC & \TBC & \TBC\\
3 & \TBC & \TBC & \TBC & \TBC & \TBC & \TBC & \TBC & \TBC\\
4 & \TBC & \TBC & \TBC & \TBC & \TBC & \TBC & \TBC & \TBC\\
5 & \TBC & \TBC & \TBC & \TBC & \TBC & \TBC & \TBC & \TBC\\
6 & \TBC & \TBC & \TBC & \TBC & \TBC & \TBC & \TBC & \TBC\\
7 & \TBC & \TBC & \TBC & \TBC & \TBC & \TBC & \TBC & \TBC\\
8 & \TBC & \TBC & \TBC & \TBC & \TBC & \TBC & \TBC & \TBC\\
9 & \TBC & \TBC & \TBC & \TBC & \TBC & \TBC & \TBC & \TBC\\
10 & \TBC & \TBC & \TBC & \TBC & \TBC & \TBC & \TBC & \TBC\\
11 & \TBC & \TBC & \TBC & \TBC & \TBC & \TBC & \TBC & \TBC\\
12 & \TBC & \TBC & \TBC & \TBC & \TBC & \TBC & \TBC & \TBC\\
13 & \TBC & \TBC & \TBC & \TBC & \TBC & \TBC & \TBC & \TBC\\
14 & \TBC & \TBC & \TBC & \TBC & \TBC & \TBC & \TBC & \TBC\\
15 & \TBC & \TBC & \TBC & \TBC & \TBC & \TBC & \TBC & \TBC\\
16 & \TBC & \TBC & \TBC & \TBC & \TBC & \TBC & \TBC & \TBC\\
17 & \TBC & \TBC & \TBC & \TBC & \TBC & \TBC & \TBC & \TBC\\
18 & \TBC & \TBC & \TBC & \TBC & \TBC & \TBC & \TBC & \TBC\\
19 & \TBC & \TBC & \TBC & \TBC & \TBC & \TBC & \TBC & \TBC\\
20 & \TBC & \TBC & \TBC & \TBC & \TBC & \TBC & \TBC & \TBC\\
\bottomrule
\end{tabular*}
\end{table*}

\begin{table}[t]
\caption{Counts derived from Table~\ref{tab:included}. All cells must be computed
from the extraction sheet, not estimated.}
\label{tab:rq3counts}
\centering
\footnotesize
\begin{tabular}{@{}p{4.7cm}r@{}}
\toprule
Category & Studies\\
\midrule
\multicolumn{2}{@{}l}{\emph{Model family} (RQ2)}\\
Random forest & \TBC\\
Support vector machine & \TBC\\
Neural network or deep learning & \TBC\\
Regression or PLS & \TBC\\
Other & \TBC\\
\midrule
\multicolumn{2}{@{}l}{\emph{Validation type} (RQ3)}\\
Independent set & \TBC\\
Spatially blocked & \TBC\\
Random split or CV only & \TBC\\
Not reported & \TBC\\
\midrule
\multicolumn{2}{@{}l}{\emph{Noise tested} (RQ3)}\\
Yes & \TBC\\
No & \TBC\\
\midrule
\multicolumn{2}{@{}l}{\emph{Quality score} (Table~\ref{tab:quality})}\\
6 (maximum) & \TBC\\
4--5 & \TBC\\
3 & \TBC\\
Below 3 ($\dagger$) & \TBC\\
\midrule
\multicolumn{2}{@{}l}{\emph{Evidence origin}}\\
SSA-demonstrated & \TBC\\
Transferable & \TBC\\
\bottomrule
\end{tabular}
\end{table}

% =====================================================================
\section{Discussion}
% =====================================================================
Laboratory analysis is the standard for soil health assessment but is out of
reach for many farmers~\cite{ref1,ref3}. Sensors and ML each address part of the
problem~\cite{ref4,ref6}, and edge AI shows that near-laboratory accuracy is
possible with portable tools~\cite{ref10,ref11}. The reviewed literature
therefore supports predicting soil properties from field sensor data with ML and
evaluating the models against laboratory results on African soils.

\noindent\textbf{Strength of the evidence.} The final evidence base is 20 studies,
from two sources and snowballing, screened by one reviewer. Findings from it show
where the literature points and do not prove that a gap exists. The effect is
larger where a finding rests on few studies or on studies scoring below 3, and
each finding in Section~\ref{sec:results} is reported with and without those
studies for that reason.

\noindent\textbf{Machine learning reliability.} For ML-based soil prediction, the
validation design affects how far an accuracy figure can be trusted. Soil samples
from nearby locations are similar, so random splits place near-duplicates in both
training and test sets and can overstate accuracy; spatially blocked validation,
which holds out whole areas, gives a harder and more realistic test of
performance on a new farm. Field sensors add a second problem, because their
readings are noisier than laboratory values, so a model that performs well on
clean data may degrade quickly on field data. These points are why the review
records the validation type (Table~\ref{tab:quality}, Q4) and any noise test for
every ML study.

\noindent\textbf{Alternative explanations for the gap.} The absence of Ugandan
work that predicts soil properties from field sensor data with ML may mean that
none exists. It may also reflect (i) work published as theses, reports or in
journals that the two sources do not index (excluded here under EC1); (ii)
different terminology, such as ``rapid soil test'' or ``soil spectroscopy'',
which the string only partly covers; (iii) language limits, since only
English-language studies were included; and (iv) timing, since low-power hardware
and applied ML are recent in this setting. The gap is therefore best stated as
``not found in the sources searched''.

\noindent\textbf{Transferability.} Results from outside SSA are labelled
transferable and should not be read as SSA evidence. Sensor readings and ML
calibrations can depend on soil type, moisture and temperature, so a model trained
elsewhere may not hold on Ugandan soils until tested there.

% =====================================================================
\section{Research Gaps}\label{sec:gaps}
% =====================================================================
% IMP6: the reviewed version drew its gap statement from Table VI, which
% Section IV had already declared background only. The counts now come from the
% included set, which is what the PRISMA figure promises.
The gaps are stated from the included studies in Table~\ref{tab:included}, and
from the counts in Table~\ref{tab:rq3counts}:

\begin{enumerate}
\item \textbf{Regional evidence.} \TBC{} of the 20 included studies are
SSA-demonstrated and \TBC{} are set in Uganda. No study predicting soil
properties from field sensor data with ML on Ugandan soils was found in the
sources searched.
\item \textbf{Validation design.} \TBC{} of the 20 included studies used a random
split or cross-validation only, so reported accuracy for those studies is not
evidence of performance on new fields.
\item \textbf{Noise robustness.} \TBC{} of the 20 included studies tested the
effect of sensor noise or input error, although field sensor readings are noisier
than the laboratory values the models are trained against.
\item \textbf{Combination.} \TBC{} studies combined field sensing, ML prediction
and a validation design that holds out ground; the proposed direction combines
all three.
\end{enumerate}

% =====================================================================
\section{Limitations}\label{sec:limits}
% =====================================================================
Only two sources were searched. Web of Science, ScienceDirect and African and
agricultural indexes such as AJOL and AGRIS were not used, so relevant
publications elsewhere may be missed; extending the search to them is the first
step for the next version. Google Scholar results were limited to the first 100
per query. Only English-language studies from 2015 onward were included, which
excludes French- and Portuguese-language work from SSA. Remote-sensing-based soil
mapping was outside the search, so studies using satellite or drone imagery may
be missed. Screening was done by a single reviewer, which raises the risk of
selection bias, and no inter-rater agreement was measured. GIS mapping and
economic analysis are outside the scope, so the review says nothing about them.
The sources found also do not describe land use by commercial farmers in detail,
so conclusions about commercial farms rest on the farm-type labels in the
included studies. Finally, the gaps reported here show limits of the evidence
found and are not proof that no study exists. The framework in
Section~\ref{sec:method} models a field sensor rather than using one, so the
sensor findings of this review define a specification for that model and do not
validate a deployed device.

% =====================================================================
\section{Methodological Approach for the Proposed Empirical Study}\label{sec:method}
% =====================================================================
The proposed study answers the validation and robustness questions directly. It
has two parts, which share one modelling pipeline: an analysis of public soil
datasets (Part B) and a simulation of field sensors (Part A). Fig.~\ref{fig:pipe}
shows the stages, and the steps below define each one.

Neither part ingests field sensor data. Part B uses laboratory measurements as
proxy variables, and Part A generates sensor readings from a measurement model
whose parameters are assumptions until they are sourced. The sensor pathway is
therefore specified and simulated rather than implemented; the hardware stage is
separate work and is not reported here.

\begin{figure}[t]
\centering
% Place the pipeline schematic beside this file as pipeline.pdf (or .png).
\includegraphics[width=\columnwidth]{pipeline.pdf}
\caption{Methodological pipeline for the proposed study.}
\label{fig:pipe}
\end{figure}

\noindent\textbf{Step 1: data and targets.}
% IMP7: the reviewed version made pH and SOC equal targets in both parts and
% called moisture a covariate. In the implementation, Part B predicts SOC with
% pH as a feature, and moisture exists only as a simulated channel in Part A.
% Part A predicts both pH and SOC. State it that way.
Part A predicts soil pH and soil organic carbon (SOC) from simulated sensor
channels, and moisture is one of those channels rather than a covariate. Part B
predicts SOC, with pH among its input features; pH is not a target in Part B
because a pH meter returns pH directly and a model is not needed for it.

Part B uses public datasets that hold coordinates and laboratory values
(Table~\ref{tab:datasets}). A dataset is selected if it has coordinates and
laboratory SOC or pH, a documented licence, and samples from Africa. Records with
missing coordinates or targets are removed, units are harmonised, duplicates are
dropped, and the dataset version is recorded. The Ugandan subset is counted and
reported only where the data are measured; a country share computed from
generated data is not reported as a subset.

\begin{table}[t]
\caption{Candidate public datasets.}
\label{tab:datasets}
\centering
\footnotesize
\begin{tabular}{@{}p{1.5cm}p{3.9cm}p{0.9cm}@{}}
\toprule
Dataset & Content & Source\\
\midrule
Open Soil Spectral Library & 91,631 mid-infrared and 65,063 visible-near-infrared scans with laboratory data from more than 60,000 locations; includes AfSIS and ICRAF-ISRIC libraries and a portable-instrument NIR subset scanned on 9 devices & \cite{ref18}\\
Africa Soil Profiles Database v1.2 & 18,532 legacy profiles for SSA; analytical data for 15,564, of which 14,197 are georeferenced; profile reliability rated 1--4 & \cite{ref19}\\
WoSIS snapshot (2019) & 196,498 georeferenced profiles from 173 countries, including organic carbon and pH & \cite{ref20}\\
% IMP8: AfSIS Phase I is the dataset actually loaded by the implementation and
% was absent from this table, so the results cited a dataset the paper never
% listed. The row is added; confirm the citation form in \cite{ref21}.
AfSIS Phase I soil data & Georeferenced field observations with laboratory wet chemistry (pH, electrical conductivity, particle-size fractions, carbon) for 19 countries; redistributed under the Open Data Commons Attribution licence & \cite{ref21}\\
\bottomrule
\end{tabular}
\end{table}

\noindent\textbf{Step 2: sensor simulation (Part A).} A synthetic field with
spatially correlated pH, SOC and moisture is surveyed at a fixed number of
points. Each ``sensor'' reading is the true value passed through a measurement
model with gain, offset, random noise, calibration drift and cross-sensitivity to
moisture. Parameters are taken from sensor datasheets or bench tests, and any
value without a source is labelled an assumption and counted as such in the
output. SOC in Part A is generated from the same latent fields that the sensors
observe, through a coupling coefficient $c$; $c$ is therefore an assumption that
is swept, not a property of soil, and Section~\ref{sec:hyp} states the
consequence for H3.

% IMP9: nested tuning is a Part B property. Part A uses default hyperparameters
% by design, and the reviewed version stated the tuning requirement for both.
\noindent\textbf{Step 3: models.} In Part B, five learners are compared: a
mean-value baseline, ridge regression, partial least squares, random forest and
gradient boosting. Hyperparameters are tuned by nested cross-validation inside
the training folds only, and random seeds are fixed. Part A compares four of
these learners at their default hyperparameters, because its purpose is to trace
the effect of sensor error rather than to maximise accuracy; Part A performs no
tuning.

\noindent\textbf{Step 4: validation.}
% IMP10: real datasets often ship a spatial hierarchy, and using it is both more
% defensible and more informative than clustering coordinates; the effective
% number of units also has to be reported, because it, not the row count, sets
% the precision of a spatially blocked estimate.
Every model is evaluated twice on the same data: with random 5-fold
cross-validation and with spatially blocked cross-validation. Blocks are formed
from the site and cluster identifiers supplied with the dataset where these
exist, and by clustering coordinates where they do not; whole groups are held
out either way. Performance is reported as $R^2$, RMSE and the ratio of
performance to interquartile distance (RPIQ) (Eq.~\ref{eq:metrics}):

\begin{equation}
\mathrm{RMSE}=\sqrt{\frac{1}{n}\sum_{i=1}^{n}(y_i-\hat{y}_i)^2},
\qquad
\mathrm{RPIQ}=\frac{\mathrm{IQR}(y)}{\mathrm{RMSE}}.
\label{eq:metrics}
\end{equation}

The validation gap is the $R^2$ under random folds minus the $R^2$ under spatial
blocks. Repeating the procedure over multiple seeds gives a spread for each
figure. Because the precision of a spatially blocked estimate depends on the
number of independent units and not on the number of rows, the number of distinct
sites, clusters and blocks used is reported with every result.

Prediction intervals are also reported. Their nominal coverage is fixed in
advance and their empirical coverage is measured on held-out blocks, so that a
quoted interval can be checked rather than assumed.

\noindent\textbf{Step 5: noise tests.} Gaussian noise is added to the inputs at
levels from 0 to 100 percent of each feature's standard deviation. Models are
tested with clean training and noisy test inputs, and with noisy training and
noisy test inputs.
% IMP11: the comparison with instrument-to-instrument variation is not part of
% the present results and must not be described as though it were. The common
% multiplier is also disclosed, because it is weaker than channel-specific error.
The same multiplier is applied to every feature in Part B, so channel-specific
error is modelled in Part A only. For Part A the noise levels follow the
simulated sensor model. Comparison with instrument-to-instrument variation in the
portable-instrument NIR subset of the Open Soil Spectral Library is planned and is
not part of the present results.

\noindent\textbf{Step 6: controls.} A model trained on shuffled targets must score
near zero $R^2$, and a mean-value predictor sets the lower baseline. The block
size used for spatial validation is varied to check that conclusions do not
depend on it.
% IMP12: the block size is not the only arbitrary choice; the block boundaries
% are too, and the implementation varies both.
Because the boundaries between blocks are as arbitrary as their size, the
partitioning is also redrawn with a different seed and the resulting spread is
reported.

\noindent\textbf{Step 7: reporting.} Code, seeds and dataset versions are released
with the results. Artefacts are bound together by hash, so a model cannot be
served against a results file it was not produced from; the manifest records the
dataset version and the model digest. Results from Part A are reported as
simulations under stated assumptions and are not claims about real sensors.
Results from Part B are reported for the subset and region actually used, and
where the data are a labelled synthetic stand-in, every reported figure says so.

% IMP13: the implementation also serves predictions, which no step described.
\noindent\textbf{Step 8: delivery.} The pipeline is exposed as a service that
accepts one sample at a time, rejects physically impossible inputs, flags inputs
outside the range seen in training, and returns the estimate together with its
interval and the provenance of the model that produced it. A prediction is
therefore traceable to a specific dataset version and model digest.

% =====================================================================
\section{Hypotheses for the Proposed Empirical Study}\label{sec:hyp}
% =====================================================================
This section is separate from the review findings. The hypotheses follow from the
research gap and the machine learning concerns above, and will be tested with the
methodology in Section~\ref{sec:method}. They are not results of this review.

\begin{itemize}
\item \textbf{H1 (validation design):} On the same dataset, SOC models score
lower under spatially blocked validation than under random $k$-fold
cross-validation. The difference is reported per seed, together with the number
of independent spatial units on which it rests.
\item \textbf{H2 (noise robustness):} Prediction accuracy falls as input noise
rises, and the rate of decline differs between model families.
% IMP14: the reviewed H3 asserted that SOC would not be predicted well from the
% sensor channels. In Part A, SOC is generated from those same channels through
% the coupling coefficient, so that clause cannot be falsified: the outcome is
% set by construction. H3 is restated so that it tests something.
\item \textbf{H3 (sensor feasibility):} For soil pH, models using simulated sensor
channels retain a spatially validated $R^2$ of at least 0.50 up to a stated error
scale, and lose it beyond that scale. For SOC, skill under spatial validation is
governed by the coupling $c$ with which SOC is generated from the same latent
fields the sensors observe; it is therefore reported for $c \in \{0, 0.5, 0.9\}$
and is not evidence about real soils. In its present state the sensor model
sources none of its parameters from datasheets or bench tests, so the stated
error scale is an assumption rather than a measured instrument property.
\end{itemize}

% =====================================================================
\section{Conclusion}
% =====================================================================
This review examined how ML has been used to predict soil properties from field
sensor data among farmers in SSA. Of the 262 records identified, 20 studies were
included. In the sources searched, no study that predicts soil properties from
field sensors with ML on Ugandan soils was found. The review treats validation
design and noise robustness as the central questions for such models. The
evidence base is small, so the next steps are to widen the search, add a second
screener, and carry out the empirical study described in
Section~\ref{sec:method}.

% IMP15: the artefact's status is stated in the paper rather than left to the
% reader to discover from the results files.
A reference implementation of the framework exists and is reported alongside the
review. Its Part B results are currently produced from a synthetic stand-in
dataset that is labelled as such in every artefact it generates; its Part A
results have not yet been produced. No part of the implementation processes field
sensor data, and neither set of results is a claim about real soils or real
instruments; the paper says so wherever the figures appear.

\section*{Acknowledgment}
The author thanks Uganda Christian University and the Department of Computer
Science for academic support during this study.
% IMP16: attribution required by the AfSIS data licence.
The AfSIS Phase I soil data used by the reference implementation are
redistributed by the World Agroforestry Centre (ICRAF) under the Open Data
Commons Attribution (ODC-By) licence, and are used here with attribution.

\begin{thebibliography}{00}
\bibitem{ref1} E.~Chlouveraki \emph{et al.}, ``Estimation of soil properties using hyperspectral imaging and machine learning,'' \emph{Smart Agric. Technol.}, vol.~10, art. no.~100790, 2025.
\bibitem{ref2} S.~Elaete, J.~Ilukor, R.~E. Isoto, A.~Moruleng, and H.~N. Opolot, ``Comparing indigenous and scientific (Palin test) soil fertility assessment approaches for better soil health monitoring, management, and farm decision-making in Uganda,'' \emph{Discover Soil}, vol.~2, art. no.~29, 2025, doi: 10.1007/s44378-025-00058-z.
\bibitem{ref3} Food and Agriculture Organization of the United Nations, Global Soil Partnership, ``Uganda rolls out the Global Soil Doctors Programme to bring soil health closer to farmers in Karamoja.'' [Online]. Available: https://www.fao.org/global-soil-partnership/resources/highlights/detail/en/c/1757615/
\bibitem{ref4} M.~R. Islam, K.~Oliullah, M.~M. Kabir, M.~Alom, and M.~F. Mridha, ``Machine learning enabled IoT system for soil nutrients monitoring and crop recommendation,'' \emph{J. Agric. Food Res.}, vol.~14, art. no.~100880, 2023, doi: 10.1016/j.jafr.2023.100880.
\bibitem{ref5} D.~Neina, ``The role of soil pH in plant nutrition and soil remediation,'' \emph{Appl. Environ. Soil Sci.}, vol.~2019, art. no.~5794869, 2019, doi: 10.1155/2019/5794869.
\bibitem{ref6} G.~N. Shwetha and G.~J. Bhat, ``Sensor based framework for soil nutrients prediction using deep learning techniques,'' \emph{Agric. Eng. Int.: CIGR J.}, vol.~27, no.~4, pp.~291--301, 2025.
\bibitem{ref7} E.~Mamabolo \emph{et al.}, ``Application of precision agriculture technologies for crop protection and soil health,'' \emph{Smart Agric. Technol.}, vol.~12, art. no.~101270, 2025.
\bibitem{ref8} S.~Getahun, H.~Kefale, and Y.~Gelaye, ``Application of precision agriculture technologies for sustainable crop production and environmental sustainability: A systematic review,'' \emph{Sci. World J.}, vol.~2024, art. no.~2126734, 2024, doi: 10.1155/2024/2126734.
\bibitem{ref9} A.~Abdullah, E.~I. Raharjo, M.~Iwan, R.~Faizal, and Maryogi, ``Development of an IoT-based soil nutrient monitoring and GIS mapping system for precision agriculture,'' \emph{Adv. Sustain. Sci., Eng. Technol.}, vol.~7, no.~4, 2025.
\bibitem{ref10} A.~Erler \emph{et al.}, ``Mobile laser-induced breakdown spectroscopy for future application in precision agriculture---A case study,'' \emph{Sensors}, vol.~23, no.~16, art. no.~7178, 2023, doi: 10.3390/s23167178.
\bibitem{ref11} R.~Kasemi, ``EdgeSoil: Low-cost soil sensor and edge AI fusion in agricultural robotics,'' in \emph{Proc. Austrian Robot. Workshop (ARW)}, Linz, Austria, 2023, pp.~1--6.
\bibitem{ref12} M.~J. Page \emph{et al.}, ``The PRISMA 2020 statement: An updated guideline for reporting systematic reviews,'' \emph{BMJ}, vol.~372, art. no.~n71, 2021, doi: 10.1136/bmj.n71.
\bibitem{ref13} B.~Kitchenham and S.~Charters, ``Guidelines for performing systematic literature reviews in software engineering,'' Keele Univ. and Univ. of Durham, EBSE Tech. Rep. EBSE-2007-01, 2007.
\bibitem{ref14} M.~L. Rethlefsen \emph{et al.}, ``PRISMA-S: An extension to the PRISMA statement for reporting literature searches in systematic reviews,'' \emph{Syst. Rev.}, vol.~10, art. no.~39, 2021, doi: 10.1186/s13643-020-01542-z.
\bibitem{ref15} Uganda Bureau of Statistics, ``Annual Agricultural Survey 2019: statistical release,'' Kampala, Uganda, 2022. [Online]. Available: https://www.ubos.org/wp-content/uploads/2022/04/Uganda\_UBOS\_StatRelease\_AAS2019-Final.pdf
\bibitem{ref16} Ministry of Finance, Planning and Economic Development, Uganda, ``Agro-industrialisation: Agriculture remains the backbone of Uganda's economy,'' Knowledge Centre summary of the 2024 National Population and Housing Census, 2025. [Online]. Available: development.finance.go.ug
\bibitem{ref17} Uganda Bureau of Statistics and Ministry of Agriculture, Animal Industry and Fisheries, ``Census of Agriculture 2008--2009,'' FAO Food and Agriculture Microdata Catalogue. [Online]. Available: https://microdata.fao.org/index.php/catalog/1594
\bibitem{ref18} J.~L. Safanelli \emph{et al.}, ``Open Soil Spectral Library (OSSL): Building reproducible soil calibration models through open development and community engagement,'' \emph{PLoS ONE}, 2025, doi: 10.1371/journal.pone.0296545.
\bibitem{ref19} ISRIC---World Soil Information, ``Africa Soil Profiles Database, version 1.2,'' ISRIC Data Hub. [Online]. Available: https://data.isric.org/geonetwork/srv/api/records/b88870b4-6af8-4e78-a3ac-38871d757525
\bibitem{ref20} ISRIC---World Soil Information, ``Standardised soil profile data for the world (WoSIS snapshot, September 2019),'' doi: 10.17027/isric-wdcsoils.20190901.
% IMP17: reference for the dataset the implementation actually loads.
% CONFIRM the citation form before submission: the data are distributed by
% ICRAF/World Agroforestry as the AfSIS Phase I soil dataset under ODC-By.
% The entry below is a placeholder in the correct style, not a verified record.
\bibitem{ref21} World Agroforestry Centre (ICRAF), ``Africa Soil Information Service (AfSIS) Phase I soil data,'' Nairobi, Kenya, 2014. [Online]. Available: https://data.worldagroforestry.org/
\end{thebibliography}

\end{document}
